% Generated draft: requires mathematical and source review.
\documentclass[11pt]{article}
\usepackage[margin=25mm]{geometry}
\usepackage{fontspec}
\setmainfont{Latin Modern Roman}
\usepackage{amsmath,amssymb,mathtools,unicode-math}
\setmathfont{Latin Modern Math}
\usepackage{enumitem,needspace,xurl,hyperref}
\setlength{\parindent}{0pt}
\setlength{\parskip}{5pt}
\setlength{\emergencystretch}{3em}
\newcommand{\N}{\mathbb{N}}
\newcommand{\Z}{\mathbb{Z}}
\newcommand{\Q}{\mathbb{Q}}
\newcommand{\R}{\mathbb{R}}
\newcommand{\C}{\mathbb{C}}
\newcommand{\F}{\mathbb{F}}
\newcommand{\A}{\mathbb{A}}
\newcommand{\PP}{\mathbb{P}}
\newcommand{\E}{\mathbb{E}}
\newcommand{\eps}{\varepsilon}
\newcommand{\dd}{\,\mathrm{d}}
\begin{document}
\section*{UnsolvedMath ID: 800104}
\textbf{Source problem code:} catalogue entry
\subsection{Definitions and mathematical statement}
Let \(G\) be a finitely generated amenable group that is orderable in the sense intended by the source. Here amenability means that there exists a finitely additive probability measure \(m\colon\mathcal P(G)\to[0,1]\) on all subsets of \(G\), invariant under left translation. The supplied record does not specify whether ``orderable'' means left-orderable or bi-orderable; these hypotheses are therefore not identified with one another here.

Let \(\operatorname{Homeo}_+(\R)\) denote the group of orientation-preserving homeomorphisms of \(\R\), and let
\[
\varphi\colon G\longrightarrow\operatorname{Homeo}_+(\R)
\]
be a group homomorphism. Assume that the action has no global fixed point, meaning
\[
\bigcap_{g\in G}\{x\in\R:\varphi(g)(x)=x\}=\varnothing,
\]
or, equivalently, \(\forall x\in\R\;\exists g\in G\) such that \(\varphi(g)(x)\ne x\). No faithfulness assumption is supplied.

An orientation-preserving affine action is a homomorphism \(\rho\colon G\to\operatorname{Aff}_+(\R)\), where
\[
\operatorname{Aff}_+(\R)=\{x\mapsto ax+b:a>0,\ b\in\R\}.
\]
Thus \(\rho(g)(x)=a_gx+b_g\), with \(a_g>0\), \(a_{gh}=a_ga_h\), and \(b_{gh}=a_gb_h+b_g\).

The source asks whether, for every such \(G\) and \(\varphi\), the following dichotomy holds:
\[
\left[\varphi\text{ is semi-conjugate to an affine action}\right]
\quad\text{or}\quad
\left[\varphi\text{ is laminar}\right].
\]
The disjunction is not asserted to be exclusive. The terms ``semi-conjugate'' and ``laminar'' are retained from the source. The supplied record does not define laminar actions or specify the direction and regularity, monotonicity, surjectivity, or properness requirements of a semi-conjugacy. Consequently this is a statement-only transcription of the proposed dichotomy; its fully precise formulation requires those source conventions, as well as the intended meaning of orderability.
\subsection{Short English statement}
For every finitely generated amenable orderable group acting on \(\R\) by orientation-preserving homeomorphisms with no common fixed point, must the action be semi-conjugate to an affine action or be laminar? The supplied record leaves the meanings of orderability, semi-conjugacy, and laminarity insufficiently specified for a complete formal statement.
\subsection{Sources}
\begin{itemize}
\item UnsolvedMath local registry, permanent ID \(800104\), exact source problem code \texttt{catalogue entry}, titled ``Foliations and Diffeomorphism Groups --- Question 45.'' The supplied registry attributes the problem to Nariman and Yazdi, proposed in 2025, and labels it \texttt{open}; this label is supplied curation, not an independently verified mathematical or literature assessment.
\item Nariman and Yazdi, \emph{Problem List on Foliations and Diffeomorphism Groups} (2025), Question 45. Supplied source URL: \url{https://arxiv.org/abs/2503.18274}. The record reports access on 2026-10-16 and extraction by \texttt{arxiv\_latex}; these provenance details are supplied and were not independently verified.
\end{itemize}
\end{document}
